\documentclass[a4paper, serif, 10pt]{moderncv}

%% moderncv
% style and color
\moderncvstyle{banking}
\moderncvcolor{black}

%% page layout/margins
\usepackage[top=2.5cm, bottom=2.5cm, left=1.5cm, right=1.5cm]{geometry}

\nopagenumbers{}

%% Babel config for French language
% ref:
% - https://latex3.github.io/babel/guides/locale-french.html
\usepackage[french]{babel}

%% fontspec
\usepackage{fontspec}

\IfFontExistsTF{Linux Libertine}{
  \setmainfont[Ligatures={TeX, Common}, Numbers=OldStyle]{Linux Libertine}
}{}
\IfFontExistsTF{Linux Libertine O}{
  \setmainfont[Ligatures={TeX, Common}, Numbers=OldStyle]{Linux Libertine O}
}{}

%% CTeX
% CJK support, used to show CJK characters in the resume
%
% - fontset=none: disable builtin fontset but instead set the CJK font manually
% - heading=false: disable ctex heading
% - punct=kaiming: use kaiming punctuations styles for CJK
% - scheme=plain: use plain scheme, do not override \normalsize font size
% - space=auto: space settings for CJK characters
%
% ref:
% - http://ctan.mirrorcatalogs.com/language/chinese/ctex/ctex.pdf
\usepackage[UTF8, heading=false, punct=kaiming, scheme=plain, space=auto]{ctex}

\IfFontExistsTF{Noto Serif CJK SC}{
  \setCJKmainfont{Noto Serif CJK SC}
}{}
\IfFontExistsTF{Noto Sans CJK SC}{
  \setCJKsansfont{Noto Sans CJK SC}
}{}

%% line spacing
\usepackage{setspace}
\setstretch{1.125}

%% URL styling - use normal text instead of monospace
\urlstyle{same}

% Auto-underline all links
% Defer until \begin{document} so hyperref is already loaded
\AtBeginDocument{%
  \let\oldhref\href
  \renewcommand{\href}[2]{\oldhref{#1}{\underline{#2}}}%
}

%% Patch \httplink and \httpslink to handle full URLs
% The original moderncv macros always prepend http:// or https:// to URLs.
% This causes issues when the URL already has a protocol (e.g., https://example.com)
% resulting in malformed URLs like https://https://example.com.
% This patch checks if the URL already contains "://" and uses it directly if so.
% Note: We use \str_set:Nx (x-expansion) to expand macros like \@homepage first.
\ExplSyntaxOn
\str_new:N \g_yamlresume_url_str

\RenewDocumentCommand{\httpslink}{O{}m}{
  \str_set:Nx \g_yamlresume_url_str {#2}
  \str_if_in:NnTF \g_yamlresume_url_str {://}
    { \url{#2} }
    { \url{https://#2} }
}

\RenewDocumentCommand{\httplink}{O{}m}{
  \str_set:Nx \g_yamlresume_url_str {#2}
  \str_if_in:NnTF \g_yamlresume_url_str {://}
    { \url{#2} }
    { \url{http://#2} }
}
\ExplSyntaxOff

\name{Camille Laurent}{}
\title{Responsable de la réussite client}
\phone[mobile]{+33 1 44 55 66 77}
\email{camille.laurent@example.fr}
\homepage{https://www.camille-laurent.fr}

\address{12 Rue de la Paix, Paris, Île-de-France, France, 75002}{}{}

\extrainfo{{\small \faLinkedin}\ \href{https://www.linkedin.com/in/camillelaurent}{@camillelaurent} {} {} {} • {} {} {} 
{\small \faTwitter}\ \href{https://twitter.com/camillelaurent}{@camillelaurent}}

\begin{document}

\maketitle

\section{Informations de base}

\cvline{}{Responsable de la réussite client avec plus de 6 ans d'expérience dans le secteur SaaS, spécialisée dans l'accompagnement des clients grands comptes.
%
\begin{itemize}
\item Expertise en gestion de portefeuille client, onboarding, et fidélisation.
\item Forte capacité à établir des relations de confiance et à identifier les opportunités de croissance.
\item Solide compréhension des enjeux techniques et métier, permettant de traduire les besoins clients en solutions concrètes.
\end{itemize}}
\section{Formation}

\cventry{sept. 2014–juin 2016}
        {Master, Management et Stratégie d'Entreprise, Score : Mention Bien}
        {Université Paris-Dauphine}
        {\url{https://www.dauphine.psl.eu/}}
        {}
        {\begin{itemize}
\item Formation spécialisée en management et stratégie, avec un focus sur la relation client et le développement commercial.
\item Projet de fin d'études sur l'optimisation de l'expérience client dans le secteur SaaS.
\end{itemize}
\textbf{Cours} : Stratégie d'entreprise, Marketing B2B, Négociation commerciale, Gestion de la relation client, Analyse financière}

\cventry{sept. 2011–juin 2014}
        {Licence, Économie et Gestion, Score : Mention Assez Bien}
        {Université Paris 1 Panthéon-Sorbonne}
        {\url{https://www.pantheonsorbonne.fr/}}
        {}
        {\begin{itemize}
\item Solide base en économie et gestion, avec une introduction aux techniques de vente et de négociation.
\end{itemize}
\textbf{Cours} : Microéconomie, Macroéconomie, Comptabilité, Statistiques, Marketing}
\section{Expérience professionnelle}

\cventry{janv. 2021–Present}
        {Responsable de la réussite client}
        {SaaSify}
        {\url{https://www.saasify.fr}}
        {}
        {\begin{itemize}
\item Gestion d'un portefeuille de 50 clients grands comptes, représentant un revenu annuel récurrent de 5M€.
\item Élaboration et mise en œuvre de plans de réussite client personnalisés, augmentant le taux de rétention de 15\%.
\item Collaboration étroite avec les équipes produit et ventes pour identifier les opportunités d'upsell et de cross-sell.
\item Animation de revues trimestrielles avec les clients pour suivre les indicateurs de performance et aligner les objectifs.
\end{itemize}
\textbf{Mots-clés} : Gestion de portefeuille, Onboarding, Fidélisation, SaaS, Négociation}

\cventry{mars 2019–déc. 2020}
        {Responsable de la réussite client}
        {CloudNext}
        {\url{https://www.cloudnext.com}}
        {}
        {\begin{itemize}
\item Suivi de 30 clients PME/ETI, avec un focus sur l'adoption produit et la satisfaction client.
\item Mise en place d'un programme d'onboarding structuré, réduisant le time-to-value de 30\%.
\item Analyse des données d'utilisation pour anticiper les risques de churn et proposer des actions correctives.
\end{itemize}
\textbf{Mots-clés} : Adoption produit, Satisfaction client, Analyse de données, Onboarding}

\cventry{sept. 2016–févr. 2019}
        {Chargé de clientèle}
        {TechSolutions}
        {\url{https://www.techsolutions.fr}}
        {}
        {\begin{itemize}
\item Gestion d'un portefeuille de 100 clients TPE/PME dans le secteur des logiciels.
\item Développement des ventes additionnelles et suivi de la satisfaction client.
\item Collaboration avec l'équipe support pour résoudre les problèmes clients dans les délais impartis.
\end{itemize}
\textbf{Mots-clés} : Vente, Relation client, Support}
\section{Langues}

\cvline{Français}{Niveau natif ou bilingue}
\cvline{Anglais}{Niveau professionnel complet \hfill \textbf{Mots-clés} : TOEIC 950, TOEFL 110}
\cvline{Espagnol}{Niveau élémentaire}
\section{Compétences}

\cvline{Gestion de la relation client}{Expert \hfill \textbf{Mots-clés} : Salesforce, HubSpot, Zendesk, Gainsight}
\cvline{Analyse de données}{Avancé \hfill \textbf{Mots-clés} : Excel, SQL, Tableau, Google Analytics}
\cvline{Gestion de projet}{Avancé \hfill \textbf{Mots-clés} : Jira, Asana, Méthodologies Agile}
\cvline{Communication}{Expert \hfill \textbf{Mots-clés} : Présentations, Négociation, Rédaction}
\section{Récompenses}

\cventry{déc. 2022}
        {SaaSify}
        {Prix de l'Excellence Client}
        {}
        {}
        {Récompense décernée pour la meilleure performance en matière de satisfaction client et de rétention sur l'année 2022.}
\section{Certificats}

\cventry{mars 2020}
        {SuccessCo}
        {Certification Customer Success Manager}
        {\url{https://www.successco.com/certification}}
        {}
        {}
\section{Publications}

\cventry{juin 2021}
        {L'importance de l'onboarding client dans le SaaS}
        {Revue du SaaS}
        {\url{https://www.revuedusaas.fr/onboarding-client}}
        {}
        {Article analysant les meilleures pratiques pour structurer un programme d'onboarding efficace et améliorer la rétention client.}
\section{Projets}

\cventry{janv. 2022–juin 2022}
        {Mise en place d'un parcours d'onboarding standardisé pour les nouveaux clients de SaaSify.}
        {Programme d'onboarding client}
        {\url{https://www.saasify.fr/onboarding}}
        {}
        {\begin{itemize}
\item Conception et déploiement d'un programme d'onboarding en 4 étapes.
\item Réduction du time-to-value de 40\% et augmentation de la satisfaction client de 20 points.
\item Formation des équipes internes à la nouvelle méthodologie.
\end{itemize}
\textbf{Mots-clés} : Onboarding, Expérience client, SaaS}
\section{Centres d'intérêt}

\cvline{Sports}{Course à pied, Natation, Yoga}
\cvline{Voyages}{Asie, Europe, Amérique du Sud}
\cvline{Lecture}{Romans, Biographies, Essais}
\section{Bénévolat}

\cventry{sept. 2017–juin 2023}
        {Mentor}
        {Association des Anciens de Dauphine}
        {\url{https://www.dauphine-alumni.org}}
        {}
        {\begin{itemize}
\item Accompagnement de jeunes diplômés dans leur insertion professionnelle.
\item Organisation d'ateliers de préparation aux entretiens et de développement du réseau.
\end{itemize}}
    
\end{document}